\documentclass[12pt]{article}
\usepackage[utf8]{inputenc}
\usepackage[T1]{fontenc}
\usepackage[margin=1in]{geometry}
\usepackage{setspace}
\usepackage{xcolor}
\usepackage[colorlinks=true, urlcolor=blue!60!black, linkcolor=blue!60!black]{hyperref}
\setlength{\parindent}{2em}
\setlength{\parskip}{1.25em}
\renewcommand{\baselinestretch}{1.0}
\title{\textit{United States v. E.C. Knight Co.}, 156 U.S. 1 (1895)}
\author{Case Brief}
\date{}
\begin{document}
\maketitle

\section{Case Metadata}
\begin{itemize}
  \item \textbf{Case Name:} \textit{United States v. E.C. Knight Co.}, 156 U.S. 1 (1895) (No.~675)
  \item \textbf{Chapter:} Chapter 3: Federal Legislative Power
  \item \textbf{Court and Dates:} Supreme Court of the United States; argued October 24, 1894; decided January 21, 1895
  \item \textbf{Lower Courts:} U.S. Circuit Court for the Eastern District of Pennsylvania (bill dismissed); U.S. Court of Appeals for the Third Circuit (affirmed); Supreme Court affirmed
  \item \textbf{Vote:} 8--1
  \item \textbf{Majority Author:} Chief Justice Melville W. Fuller (joined by Justices Field, Gray, Brewer, Brown, Shiras, Jackson, and White)
  \item \textbf{Dissenting Opinion:} Justice John Marshall Harlan (sole dissent)
  \item \textbf{Primary Texts:} \href{https://supreme.justia.com/cases/federal/us/156/1/}{Justia U.S. Supreme Court Center}; \href{https://www.law.cornell.edu/supremecourt/text/156/1}{Cornell LII Full Opinion}
  \item \textbf{Constitutional and Statutory Provisions:} U.S. Const. art.~I, \S~8, cl.~3 (Commerce Clause); art.~I, \S~8, cl.~18 (Necessary and Proper Clause); Tenth Amendment; Sherman Antitrust Act of July 2, 1890, 26 Stat. 209
\end{itemize}

\section{Facts and Procedural History}
In March 1892, the American Sugar Refining Company---a New Jersey corporation already controlling the vast majority of sugar refining in the United States---contracted to purchase the stock of four competing Philadelphia refineries: the E.C. Knight Company, the Franklin Sugar Refining Company, Spreckels' Sugar Refining Company, and the Delaware Sugar House (\href{https://supreme.justia.com/cases/federal/us/156/1/}{156 U.S. at 3--4}). With those four acquisitions, only one independent competitor remained (the Revere Sugar Refinery of Boston, producing roughly 2\% of national output), giving the ``Sugar Trust'' control of approximately 98\% of all refined sugar manufactured in the United States (\href{https://www.law.cornell.edu/supremecourt/text/156/1}{id. at 3--4, 18}).

The United States filed a bill in equity in the U.S. Circuit Court for the Eastern District of Pennsylvania under the Sherman Antitrust Act of 1890, alleging that the stock acquisitions constituted a contract, combination, or conspiracy to restrain and monopolize interstate and foreign trade in refined sugar and seeking to cancel the agreements. The Circuit Court dismissed the bill, concluding that the contracts related only to manufacturing within Pennsylvania rather than interstate commerce (\href{https://www.law.cornell.edu/supremecourt/text/156/1}{156 U.S. at 5}). The Third Circuit affirmed, and the United States appealed to the Supreme Court.

\section{Issue(s) Presented}
Whether Congress's power under the Commerce Clause (and its exercise in the Sherman Antitrust Act of 1890) authorizes the federal government to dismantle a near-complete monopoly in the \textit{manufacture} of a nationally distributed commodity (refined sugar) where the acquired refineries are located within a single state.

\section{Holding}
No. In an 8--1 decision authored by Chief Justice Fuller, the Supreme Court affirmed dismissal of the government's antitrust suit. The Court held that the Sherman Act could not reach a monopoly in sugar manufacturing because ``commerce succeeds to manufacture, and is not a part of it'' (\href{https://supreme.justia.com/cases/federal/us/156/1/}{156 U.S. at 12}). Because manufacturing precedes interstate distribution and affects interstate commerce only ``incidentally and indirectly,'' regulating manufacturing monopolies belonged exclusively to state police power rather than Congress's Commerce Clause authority (\href{https://supreme.justia.com/cases/federal/us/156/1/}{id. at 12--17}).

\section{Reasoning}
Chief Justice Fuller built the majority's dual-federalism holding on four formalist distinctions:

\begin{enumerate}
  \item \textbf{Categorical Separation of Manufacture from Commerce:} Drawing on \textit{Kidd v. Pearson}, 128 U.S. 1 (1888), Fuller insisted that ``manufacture is transformation---the fashioning of raw materials into a change of form for use,'' whereas commerce involves buying, selling, and interstate transportation. Even if every pound of sugar refined in Philadelphia was destined for interstate shipment, ``commerce succeeds to manufacture, and is not a part of it'' (\href{https://www.law.cornell.edu/supremecourt/text/156/1}{156 U.S. at 12}).
  \item \textbf{Exclusive State Police Power vs. Federal Commerce Power:} Treating state and federal spheres as mutually exclusive, Fuller declared that ``[t]hat which belongs to commerce is within the jurisdiction of the United States, but that which does not belong to commerce is within the jurisdiction of the police power of the state'' (\href{https://www.law.cornell.edu/supremecourt/text/156/1}{156 U.S. at 12}). Suppressing a local manufacturing monopoly was therefore a matter for Pennsylvania's police power, not federal antitrust enforcement.
  \item \textbf{Direct vs. Indirect Effects:} Fuller conceded that controlling 98\% of sugar manufacturing ``involves, in a certain sense, the control of its disposition,'' but characterized that effect as ``secondary, and not the primary, sense,'' concluding that the combination affected interstate commerce ``only incidentally and indirectly'' (\href{https://supreme.justia.com/cases/federal/us/156/1/}{156 U.S. at 12, 17}). Because the government's proofs showed an acquisition of manufacturing stock rather than a direct restraint on interstate sales or transportation, the Sherman Act did not apply.
  \item \textbf{Slippery-Slope Federalism Warning:} Fuller warned that if Congress could regulate local manufacturing whenever a commodity entered interstate trade, ``comparatively little of business operations and affairs would be left for state control,'' sweeping agriculture, mining, and every branch of domestic production under federal supervision (\href{https://www.law.cornell.edu/supremecourt/text/156/1}{156 U.S. at 16}).
\end{enumerate}

\section{Concurring and Dissenting Opinions}
\textbf{Justice John Marshall Harlan (dissenting):} In a prophetic solo dissent, Justice Harlan rejected the majority's rigid boundary between production and trade. Quoting \textit{Gibbons v. Ogden} and \textit{McCulloch v. Maryland}, Harlan argued that interstate commerce is not mere transportation; it encompasses ``the purchase and sale of articles that are intended to be transported from one State to another---every species of commercial intercourse among the States'' (\href{https://www.law.cornell.edu/constitution-conan/article-1/section-8/clause-3/commerce-among-the-several-states}{156 U.S. at 22}). A trust controlling 98\% of a household necessity refined specifically for nationwide sale directly obstructed interstate commerce and exacted higher prices from consumers across every state (\href{https://www.law.cornell.edu/constitution-conan/article-1/section-8/clause-3/commerce-among-the-several-states}{id. at 33}). Because no single state could effectively police a multistate corporate combination chartered in New Jersey and operating nationwide, Harlan concluded that the Commerce Clause and Necessary and Proper Clause empowered Congress to dismantle the monopoly; stripping the national government of that power left the public helpless against interstate trusts.

\section{Key Quotations}
\begin{quote}
``Commerce succeeds to manufacture, and is not a part of it. \ldots That which belongs to commerce is within the jurisdiction of the United States, but that which does not belong to commerce is within the jurisdiction of the police power of the state.'' --- Chief Justice Fuller, \href{https://supreme.justia.com/cases/federal/us/156/1/}{156 U.S. at 12}
\end{quote}

\begin{quote}
``Doubtless the power to control the manufacture of a given thing involves, in a certain sense, the control of its disposition, but this is a secondary, and not the primary, sense, and although the exercise of that power may result in bringing the operation of commerce into play, it does not control it, and affects it only incidentally and indirectly.'' --- Chief Justice Fuller, \href{https://supreme.justia.com/cases/federal/us/156/1/}{156 U.S. at 12}
\end{quote}

\begin{quote}
``Nevertheless it does not follow that an attempt to monopolize, or the actual monopoly of, the manufacture was an attempt, whether executory or consummated, to monopolize commerce, even though, in order to dispose of the product, the instrumentality of commerce was necessarily invoked.'' --- Chief Justice Fuller, \href{https://supreme.justia.com/cases/federal/us/156/1/}{156 U.S. at 17}
\end{quote}

\begin{quote}
``Interstate commerce does not, therefore, consist in transportation simply. It includes the purchase and sale of articles that are intended to be transported from one State to another---every species of commercial intercourse among the States and with foreign nations.'' --- Justice Harlan (dissenting), \href{https://www.law.cornell.edu/constitution-conan/article-1/section-8/clause-3/commerce-among-the-several-states}{156 U.S. at 22}
\end{quote}

\section{Doctrinal Significance and Modern Legacy (Through 2026)}
\textit{E.C. Knight} marked the high-water mark of late-nineteenth-century formalist dual federalism, temporarily hobbling Sherman Act enforcement against industrial trusts and setting the template for the Court's pre-1937 direct/indirect Commerce Clause limits (e.g., \textit{Hammer v. Dagenhart} and \textit{Carter v. Carter Coal Co.}). Yet its specific antitrust holding began eroding within a decade and was thoroughly repudiated during the New Deal constitutional revolution:

\begin{itemize}
  \item \textbf{Early Erosion and Complete Repudiation:} Just ten years later, in \href{https://supreme.justia.com/cases/federal/us/196/375/}{\textit{Swift \& Co. v. United States}, 196 U.S. 375 (1905)}, Justice Holmes articulated the ``current of commerce'' doctrine to uphold a Sherman Act injunction against meatpackers, distinguishing \textit{E.C. Knight} as ``accidental, secondary, or remote.'' In \href{https://supreme.justia.com/cases/federal/us/221/1/}{\textit{Standard Oil Co. of New Jersey v. United States}, 221 U.S. 1 (1911)}, the Court broke up a massive refining trust under the Sherman Act. The Court formally buried \textit{E.C. Knight}'s manufacture/commerce wall in \href{https://supreme.justia.com/cases/federal/us/301/1/}{\textit{NLRB v. Jones \& Laughlin Steel Corp.}, 301 U.S. 1 (1937)} (upholding federal labor regulation of steel manufacturing because industrial strife had a ``close and substantial relation'' to interstate commerce) and \href{https://supreme.justia.com/cases/federal/us/317/111/}{\textit{Wickard v. Filburn}, 317 U.S. 111 (1942)} (abandoning ``mechanical'' direct/indirect labels in favor of aggregate substantial economic effects, even for home-grown wheat). Today, the Sherman Act extends to the full outer reach of Congress's commerce power (\href{https://www.law.cornell.edu/constitution-conan/article-1/section-8/clause-3/sherman-antitrust-act-of-1890-and-sugar-trust-case}{Constitution Annotated}).
  \item \textbf{Enduring Federalism Echoes Through 2026:} Although \textit{E.C. Knight}'s categorical exclusion of manufacturing is dead letter, Chief Justice Fuller's underlying structural warning---that an unlimited Commerce Clause would erase the distinction between national and local power---resurfaced in the Rehnquist and Roberts Courts' federalism revival. In \href{https://supreme.justia.com/cases/federal/us/514/549/case.pdf}{\textit{United States v. Lopez}, 514 U.S. 549 (1995)}, and \href{https://supreme.justia.com/cases/federal/us/529/598/case.pdf}{\textit{United States v. Morrison}, 529 U.S. 598 (2000)}, the Court refused to ``pile inference upon inference'' to regulate intrastate \textit{noneconomic} conduct (gun possession near schools; gender-motivated violence), echoing Fuller's insistence that Congress possesses no general federal police power. Conversely, in \href{https://supreme.justia.com/cases/federal/us/545/1/}{\textit{Gonzales v. Raich}, 545 U.S. 1 (2005)}, the Court reaffirmed \textit{Wickard} for intrastate economic commodities, with \href{https://supreme.justia.com/cases/federal/us/545/03-1454/concur.pdf}{Justice Scalia's concurrence} citing Justice Harlan's \textit{E.C. Knight} dissent to ground federal reach over intrastate commerce in the Necessary and Proper Clause. Through 2026, \textit{E.C. Knight} (alongside \href{https://supreme.justia.com/cases/federal/us/567/519/}{\textit{NFIB v. Sebelius}, 567 U.S. 519 (2012)}) remains a foundational teaching foil illustrating both the collapse of formalist production/commerce categories and the ongoing judicial struggle to define enforceable outer limits on federal power.
\end{itemize}

\section{Pedagogical Critiques and Exam Hypotheticals}
\textbf{Pedagogical Critique:} \textit{E.C. Knight} is a textbook study in the pitfalls of conceptual formalism divorced from economic reality. As Justice Harlan demonstrated, monopolizing 98\% of a nationally consumed commodity's supply is functionally indistinguishable from restraining its interstate sale and price. Yet modern students should also ask why the Court did not simply abandon all limits after 1937: \textit{E.C. Knight}'s failure forced later Courts (\textit{Lopez}, \textit{Morrison}, \textit{NFIB}) to search for a more workable boundary---such as the economic/noneconomic and activity/inactivity distinctions---rather than resurrecting the unworkable manufacture/commerce divide.

\textbf{Exam Hypothetical (Standard Contrast---1895 vs. Modern Doctrine):} A Delaware holding company acquires four intrastate semiconductor fabrication plants in Ohio, giving it 95\% of U.S. production of a chip shipped to automakers nationwide. The Department of Justice sues under the Sherman Act to unwind the merger. \textit{Analysis:} Under \textit{E.C. Knight}'s 1895 framework, the defendants would argue that chip fabrication is ``manufacture,'' not ``commerce,'' and affects interstate trade only indirectly. Under modern doctrine (\textit{Jones \& Laughlin}, \textit{Wickard}, \textit{Raich}), the government easily prevails: chip manufacturing is quintessential economic activity whose monopolization substantially affects interstate commerce.

\textbf{``Opposite Side'' Contrast (Where \textit{E.C. Knight}'s Federalism Concern Still Wins Today):} Suppose instead that Congress enacts a federal statute criminalizing purely local, noncommercial possession of a homemade tool within 500 feet of a community park, with no jurisdictional nexus to interstate commerce and no broader regulation of an interstate commercial market, arguing that local violence indirectly reduces interstate tourism. \textit{Contrast:} Here, the federalism principle animating \textit{E.C. Knight} prevails in modern dress: under \textit{Lopez} and \textit{Morrison}, Congress cannot aggregate attenuated, indirect effects of \textit{noneconomic} intrastate conduct without converting the Commerce Clause into a general police power reserved to the states.

\end{document}
